\section{Quantiles and medians}

Expectation and variance summarize a distribution by numerical averages. Another
important way to summarize a continuous distribution is to ask where a given
amount of probability has accumulated. This leads to quantiles.

The cumulative distribution function is the natural tool here. Since
\[
F_X(x)=P(X\leq x),
\]
the value \(F_X(x)\) tells us how much probability lies to the left of \(x\).
A quantile reverses this question.

\begin{definitionbox}
Let \(X\) be a random variable with CDF \(F_X\). A number \(q_p\) is called a
\textbf{\(p\)-quantile} of \(X\), where \(0<p<1\), if
\[
P(X\leq q_p)\geq p
\]
and
\[
P(X\geq q_p)\geq 1-p.
\]
When the CDF is continuous and strictly increasing, this is equivalent to
\[
F_X(q_p)=p.
\]
\end{definitionbox}

The second formulation, \(F_X(q_p)=p\), is often the one used in calculations.
The more careful definition is included because not every distribution has a
strictly increasing CDF. In this chapter, most examples are simple continuous
models where the equation \(F_X(q_p)=p\) gives the intended quantile.

\subsection*{Median}

The most important quantile is the median. A median is a point that splits the
distribution into two parts of probability at least one half on each side. For a
continuous distribution with strictly increasing CDF, the median \(m\) satisfies
\[
F_X(m)=\frac12.
\]
Equivalently,
\[
P(X\leq m)=\frac12.
\]

The median is not the same concept as the expectation. The expectation is a
center of balance. The median is a point of probability balance: half of the
probability lies to the left and half to the right.

\subsection*{Example: uniform distribution on \([0,1]\)}

For \(X\sim\operatorname{Unif}(0,1)\), the CDF is
\[
F_X(x)=x,
\]
for \(0\leq x\leq1\). The \(p\)-quantile satisfies
\[
q_p=p.
\]
For example, the median is
\[
q_{0.5}=0.5,
\]
and the first and third quartiles are
\[
q_{0.25}=0.25,
\qquad
q_{0.75}=0.75.
\]
This agrees with the interpretation of probability as length.

\subsection*{Example: triangular density}

Let
\[
f_X(x)=
\begin{cases}
2x, & 0\leq x\leq1,\\
0, & \text{otherwise}.
\end{cases}
\]
The CDF on \([0,1]\) is
\[
F_X(x)=x^2.
\]
The \(p\)-quantile solves
\[
x^2=p.
\]
Since \(x\geq0\),
\[
q_p=\sqrt p.
\]
The median is therefore
\[
q_{0.5}=\sqrt{0.5}=\frac1{\sqrt2}\approx0.707.
\]
This median is greater than \(0.5\). That makes sense because the density is
larger near \(1\), so probability accumulates more slowly at first and more
quickly later.

\subsection*{Expectation versus median}

For symmetric distributions, expectation and median often coincide. For skewed
distributions, they may differ.

In the triangular example with density \(f_X(x)=2x\) on \([0,1]\), we found
\[
\E[X]=\frac23\approx0.667,
\]
while the median is
\[
\frac1{\sqrt2}\approx0.707.
\]
These numbers are close but not equal. They summarize the distribution in
different ways. The expectation balances the density as a physical mass would be
balanced. The median divides probability into two equal halves.

\subsection*{Quartiles and interquartile range}

The quartiles are quantiles corresponding to probabilities \(0.25\), \(0.5\), and
\(0.75\). The first quartile \(Q_1\) satisfies approximately
\[
F_X(Q_1)=0.25,
\]
the median \(Q_2\) satisfies
\[
F_X(Q_2)=0.5,
\]
and the third quartile \(Q_3\) satisfies
\[
F_X(Q_3)=0.75.
\]
The interquartile range is
\[
IQR=Q_3-Q_1.
\]
It measures the length of the central interval containing about half of the
probability.

These ideas will return when empirical distributions and descriptive statistics
are studied. There, quantiles are computed from data. Here, quantiles are
computed from a theoretical distribution.

\subsection*{Reading quantiles as probability statements}

A quantile should always be interpreted probabilistically. If \(q_{0.9}=12\),
then the statement means that approximately \(90\%\) of the probability lies at
or below \(12\), and approximately \(10\%\) lies above \(12\), assuming a
continuous model with no ambiguity at the quantile.

It does not mean that \(12\) is likely as an exact value. In a continuous model,
\(P(X=12)=0\) when \(X\) has a density. The quantile is important because of the
area accumulated up to it.

\begin{summarybox}
Quantiles invert the cumulative distribution function. For continuous models
with strictly increasing CDF,
\[
F_X(q_p)=p.
\]
The median corresponds to \(p=0.5\). Quantiles describe probability positions,
not point probabilities.
\end{summarybox}
